\subsection{RADICAL OF A LIE GROUP}

PROPOSITION 23. Let G be a finite-dimensional real or complex Lie group, r the radical of L(G) (Chapter I, § 5, Definition 2) and n the largest nilpotent ideal of L(G) (Chapter I, § 4, no. 4). Let R (resp. N) be the integral subgroup of G with Lie algebra r (resp. n). Then R (resp. N) is a solvable (resp. nilpotent) Lie subgroup of G, which is invariant under every continuous automorphism of G. Every solvable (resp. nilpotent) connected normal subgroup of G is contained in R (resp. N).

The group R is solvable (§ 6, Proposition 19). Suppose that K = R. Let G' be a connected solvable normal subgroup of G. Then \(\overline{G'}\) is a connected solvable (no. 1, Corollary 2 to Proposition 1) normal Lie subgroup of G (§ 8, no. 2, Theorem 2). Hence L \((\overline{G'})\) is a solvable ideal of L(G), whence L \((\overline{G'}) \subset r\) and \(\overline{G'} \subset R\) . In particular, \(\overline{R} \subset R\) , hence R is closed and therefore is a Lie subgroup of G. Suppose that K = C. Let H be the underlying real Lie subgroup of G. If \(r'\) is the radical of L(H), \(i'r\) is a solvable ideal of L(H), whence \(r' = i'r'\) ; hence \(r \subset r' \subset r\) and, by the above, R is closed in H and therefore in G; thus R is a Lie subgroup of G. Every connected solvable normal subgroup of G is a connected solvable normal subgroup of H and therefore contained in R. Hence we have proved for K = C and for K = R that R is the largest connected solvable normal subgroup of G; therefore R is invariant under every continuous automorphism of G. The proof for N is completely analogous.

DEFINITION 1. Let G be a finite-dimensional real or complex Lie group. The radical of G is the largest connected solvable normal subgroup of G.

Remark. Even if G is connected there may be solvable normal subgroups of G not contained in the radical of G.

PROPOSITION 24. Suppose that K = R or C. Let \(G_{1}\) , \(G_{2}\) be two finite-dimensional connected Lie groups, \(R_{1}\) and \(R_{2}\) their radicals and \(\phi\) a surjective morphism of \(G_{1}\) into \(G_{2}\) . Then \(\phi(R_{1}) = R_{2}\) .

By § 3, no. 8, Proposition 28, \(\mathrm{L}(\phi)\) is surjective. Hence \(\mathrm{L}(\phi)(\mathrm{L}(\mathbb{R}_1)) = \mathrm{L}(\mathbb{R}_2)\) (Chapter I, § 6, Corollary 3 to Proposition 2). Let \(i\) be the canonical injection of \(\mathbb{R}_1\) into \(\mathbb{G}_1\) . Then the image of \(\phi \circ i\) is \(\mathbb{R}_2\) (§ 6, no. 2, Corollary I to Proposition 1).

PROPOSITION 25. Suppose that \(\mathbf{K} = \mathbf{R}\) or \(\mathbf{C}\) . Let \(\mathbf{G}_1, \mathbf{G}_2\) be finite-dimensional connected Lie groups and \(\mathbf{R}_1\) and \(\mathbf{R}_2\) their radicals. The radical of \(\mathbf{G}_1 \times \mathbf{G}_2\) is \(\mathbf{R}_1 \times \mathbf{R}_2\) . This follows from Chapter I, § 5, Proposition 4.

